\documentclass[11pt]{article}
\usepackage[a4paper,margin=22mm]{geometry}
\usepackage{fontspec}
\setmainfont{LinLibertine_R.otf}[BoldFont=LinLibertine_RB.otf,ItalicFont=LinLibertine_RI.otf,BoldItalicFont=LinLibertine_RBI.otf]
\usepackage{xcolor}
\usepackage{hyperref}
\hypersetup{colorlinks=true,urlcolor=blue}
\setlength{\parindent}{0pt}
\setlength{\parskip}{6pt}
\sloppy
\emergencystretch=2em
\begin{document}
\section*{Historical Britannica: extracted source}
Supplement to the fourth, fifth and sixth editions of the Encyclopaedia Britannica - Volume 1 - EB.14\par
{\small Source row 2281. Parser: nvidia/nemotron-3.5-parse.\par}
{\small This is model-extracted text, not a corrected transcription. OCR errors, reading order and region alignment require human review.\par}
\hrule\medskip
{\small\itshape 364\par}

{\scriptsize\color{gray} Region region-0; provider order 0.\par}
\medskip

\subsection*{A N G L E.}

{\scriptsize\color{gray} Region region-1; provider order 1.\par}
\medskip

Triscion of Elementary Geometry. They soon perceived that such problems cannot be solved by any combination of mere straight lines or circles. To this conclusion they were held directly by the application of \_Gometricical Analysis\_, a beautiful instrument of discovery which. Place had recently invented or improved. Their investigations carried some curves of a higher order than the circle, and opened to them a wide and interesting field of research.
\par

{\scriptsize\color{gray} Region region-2; provider order 2.\par}
\medskip

The analysis of the triscion of an angle, conducted in two different ways, triminates in the construction of the \_Conchoid\_, a complex curve which was first proposed by Mercedes. As the subject is very curious, and throws great light on the theory of angular magnitude, we shall here not only give both the angent methods of investigation, but subjoin a third which is due to the aspects of Newton.
\par

{\scriptsize\color{gray} Region region-3; provider order 3.\par}
\medskip

Let it be required to triest the angle BAC or the arc BC. Suppose the thing already done, and the angle BAD to be the third part of the given angle. From the point CE draw CE parallel to AD, meeting the extended diameter in \_E\_, and cutting the cir-
\par

{\scriptsize\color{gray} Region region-4; provider order 4.\par}
\medskip

C
E
D
B
E
G
A
\par

{\scriptsize\color{gray} Region region-5; provider order 5.\par}
\medskip

cumference of the circle in the point F; join FD and FA. It is obvious, that the angle BAD is the half of DAC, the remaining part of the whole angle, and therefore, equal to the angle DFG at the circumference. But AD and EC being parallel, the angle BAD is equal to AEF, which is hence equal to DFG; and subsequently FD is parallel to EB. Wherefore, the are BD is equal to GF, and the angle BAD equal to GAF. The angle AEF is thus equal to EAF, and hence the side EF is equal to AF, the radius of the circle.
\par

{\scriptsize\color{gray} Region region-6; provider order 6.\par}
\medskip

To solve the problem, therefore, it would be requisite to inflect from 0 at straight line CFE, such that the position PFB intercepted between the circumference and the diameter or its extension, should be equal to the radius of the circle. The name AD drawn parallel to this indicates line GE would cut off an angle BAD, which is the third part of the given angle BAC.
\par

{\scriptsize\color{gray} Region region-7; provider order 7.\par}
\medskip

But it is clear, that Elementary Geometry will not furnish the means of infecting CE, according to the required conditions. This must be done either tambiety that is, by repeated trials, or by the appliestion of a curve, so constituted, that every straight line drawn from the \_pide\_ C, to the \_divorteric\_ BG, shall have the portion BE, interpreted by the curve equal to AB. This curve is, from its general shape or resemblance to a \_cochal\_ shell, named the \_Conchoid\_s consist of two branches, one shows the directrix called the \_interior conchoid\_, and the other
\par

{\scriptsize\color{gray} Region region-8; provider order 8.\par}
\medskip

below it, called the \_extterior conchoid\_. The conchoid being described, will, by its interaction with the occurrence of the circle, give the point E, and consequently the position of trisecting line AD. But such a complex curve must cut the circumference in more points than one, and consequently the problem of angular friction, viewed in its generality, admits to several answers. In fact, there are always three distinct positions of the infected line CE, which will fulfill the conditions of the problem.
\par

{\scriptsize\color{gray} Region region-9; provider order 9.\par}
\medskip

It is curious to examine these different positions of the infected line. Draw AD parallel to the second iteration CE/and the DPF, AD’, and AF. Because AE is equal to E/F, the angle E/AF is equal to AEF/s, and, consequently, the exterior angle AFL is the double of either of these. But CAF being an isoceles triangle, AFL is equal to ACV, which again is equal to the alternate angle CAD’s where, CAD’s as the double of the angle AEF/s; and single likethe site the double of CFD is the circumference, the angles CFD’d and AE/F are equal; and, consequently, if ZF and EA parallel. Now the angle CAD being double of E/AF or D’AG, and the angle CAD double of DAB, the are MCD’s double of the arcs D’G and DB, which serve to com
\par

{\scriptsize\color{gray} Region region-10; provider order 10.\par}
\medskip

C
F
F
D
E
G
A
E
B
E
D'
\par

{\scriptsize\color{gray} Region region-11; provider order 11.\par}
\medskip

plete the semicircumference; where, this are DCD’ is two-third parts of the semicircumference, or can- third of the whole circumference.
\par

{\scriptsize\color{gray} Region region-12; provider order 12.\par}
\medskip

In the third position AD of the trisecting line, draw CPF parallel to A and join AF and DFT. The isoscales triangles D’AF and AFF? have equal vertical angles (and, consequently, the angles at their base are likewise equal); wheretors, AFD’s is equal to the alternatingle PFAF, and the chord D’t parallel to the diameter BG. But the reverse angle CAD’s standing on the arc CDD’s, is double of the angle CFD’s at the circumference, and, therefore, double of BFW or of BAIFs, and the angle CAD being by the first construction, likewise double of BAD, the reverse angle CAD’y together with CAD must double of BAD’ and BAD, or the arc CDD’d is double of DBD’, which completes the circumference. Hence the arc DDD’s is two-thirds of the circumference.
\par

{\scriptsize\color{gray} Region region-13; provider order 13.\par}
\medskip

It thus appears that the construction of the problem assumes three different aspects, and that the intersecting lines, to which a close analogy conducts us, mutually divide the whole circuit into equal portions. These results are perfectly conformable with the theory of angular magnitude. For if A denote the arc BC, and C the whole circumference, this arc
\par

{\scriptsize\color{gray} Region region-14; provider order 14.\par}
\medskip

{\small\itshape 10\par}

{\scriptsize\color{gray} Region region-15; provider order 15.\par}
\medskip
\newpage\section*{Source record}
bri-2281-ff9c4373e498\par
Image SHA-256: \texttt{ff9c4373e498cf1488dade04b3ed25c7}\par\texttt{d5c4da4a09a1e102e863f69b1e53fdca}\par
Dataset revision: \texttt{309f765b1ecc83e0a7f7}\texttt{762e44b4a927a7555443}\par
National Library of Scotland, Encyclopaedia Britannica Illustrated (CC0).\par
\url{https://huggingface.co/datasets/NationalLibraryOfScotland/encyclopaedia_britannica_illustrated/viewer/default/train?row=2281}
\end{document}
